\chapter*{Abstract}

Network dataplane construction faces two opposing pressures. On one side, latency and throughput targets met only by eliminating allocation from the hot path: no heap traffic, no per-packet allocation, no copy that can be fused away. On the other, a demand for modular construction: small components --- parsers, filters, encoders, ring buffers, transports --- composed into larger pipelines without surrendering the performance guarantee. The usual answer to compositionality, the abstraction layer, is itself the most common source of allocation; a dataplane that must be allocation-free by construction has to be optimized by composition, not in spite of it.

This thesis develops the algebraic theory behind it. It defines \textbf{Mol}, the category of stateful transformations, whose objects are types and whose morphisms are deterministic Mealy machines: equivalence classes of pairs $(S, \mathit{step})$ with $\mathit{step} : S \times A \to B \times S$, quotiented by the state-space bijection condition so that only observable behavior identifies a machine. Mol is symmetric monoidal under the tuple product, unit $()$; sequential composition is $\circ$ and parallel composition is $\otimes$. Its pure subcategory is equivalent to the category of types and total functions, and its effectful subcategory is isomorphic to the Kleisli category of the state monad; every hybrid molecule decomposes uniquely as $q \circ e \circ p$.

The thesis statement is that Mol is a complete algebraic foundation for zero-allocation dataplane construction: every dataplane artifact is a molecule, and its engineering properties are theorems. The catalog NT1--NT55, each with a complete proof, covers: the structure of Mol (NT1--NT8); Lawvere theories for rings, parsers, and transports, and their rewrite theories (NT9--NT19, NT54); free molecules and normal forms as a precomputable lookup table (NT20--NT24); cost enrichment with global optimality by shortest paths (NT25--NT32, NT53); Kleisli decomposition and traced feedback with Banach and Dobrushin bounds (NT33--NT45); operadic batching with syscall amortization (NT46--NT48); cut elimination as pipeline fusion, linearity as a zero-allocation guarantee, and affine typing (NT49--NT50, NT55); and situations as finite, decidable optimization problems (NT51--NT52). Classical results --- coherence, the Kleisli construction, Nerode--Hopcroft minimization, contraction theory, traced structure --- are restated and proved inside Mol; statements original to this thesis are tagged [N], classical statements [C].

The upshot is a theory in which normalization is sound and cost-preserving, minimal machines are computable, loops are provably bounded, and zero allocation is a consequence of typing rather than a performance accident. Computational claims are verified by executable Rust tools; no conjecture is presented as a theorem. A companion software standard turns the same theorems into enforceable engineering policy.
